% p161_1; source PDF page 161; manually verified editable table.
\begin{tikzpicture}[x=1cm,y=-1cm,inner sep=0pt,outer sep=0pt]
\path[use as bounding box] (0,0) rectangle (16.000,4.670);
\draw[line width=.4pt] (0.000,0.000) rectangle (8.300,0.855);
\node[align=center,text width=8.040cm,font={\linespread{1}\fontsize{10}{13.50}\selectfont\bfseries}] at (4.150,0.427) {结合类型};
\draw[line width=.4pt] (8.300,0.000) rectangle (16.000,0.855);
\node[align=center,text width=7.440cm,font={\linespread{1}\fontsize{10}{13.50}\selectfont\bfseries}] at (12.150,0.427) {处理结论};
\draw[line width=.4pt] (0.000,0.855) rectangle (8.300,1.960);
\node[align=left,text width=8.040cm,font={\linespread{1}\fontsize{10}{13.50}\selectfont}] at (4.150,1.407) {有身份者是身份犯的实行犯+无身份者是共犯};
\draw[line width=.4pt] (8.300,0.855) rectangle (16.000,1.960);
\node[align=left,text width=7.440cm,font={\linespread{1}\fontsize{10}{13.50}\selectfont}] at (12.150,1.407) {按实行犯来定罪（也即按有身份者定罪）（例 1）};
\draw[line width=.4pt] (0.000,1.960) rectangle (8.300,3.065);
\node[align=left,text width=8.040cm,font={\linespread{1}\fontsize{10}{13.50}\selectfont}] at (4.150,2.513) {无身份者是非身份犯的实行犯+有身份者是共犯};
\draw[line width=.4pt] (8.300,1.960) rectangle (16.000,3.065);
\node[align=left,text width=7.440cm,font={\linespread{1}\fontsize{10}{13.50}\selectfont}] at (12.150,2.513) {按实行犯来定罪（也即按无身份者定罪）（例 2）};
\draw[line width=.4pt] (0.000,3.065) rectangle (8.300,4.670);
\node[align=left,text width=8.040cm,font={\linespread{1}\fontsize{10}{13.50}\selectfont}] at (4.150,3.867) {A 身份者是 A 身份犯的实行犯+B 身份者是 B 身份犯的实行犯};
\draw[line width=.4pt] (8.300,3.065) rectangle (16.000,4.670);
\node[align=left,text width=7.440cm,font={\linespread{1}\fontsize{10}{13.50}\selectfont}] at (12.150,3.867) {按实行犯来定罪，存在想象竞合，择一重罪论处（例 3）};
\end{tikzpicture}
